%% ============================================================
%%  SECCIÓN 11 — Análisis de Gráficas
%%  Responsable: Minilux
%% ============================================================

\section{Análisis de Gráficas}
\label{sec:analisis_graficas}

\subsection{Gráfica 1 --- \texorpdfstring{$\gamma_i$}{gamma} vs.\ medición (primer método)}

La Fig.~\ref{fig:grafica_gamma_m1} muestra que los cinco valores de
$\gamma_i$ se distribuyen de forma dispersa en torno al promedio
$\bar{\gamma}_1 = \qty{39,73}{\milli\newton\per\meter}$, sin una
tendencia monótona creciente ni decreciente entre mediciones sucesivas.
Esto indica que las variaciones entre mediciones no responden a un
efecto sistemático acumulativo, sino a la resolución discreta de las
ranuras del brazo mayor y a la incertidumbre en la masa de los
jinetillos, ambas fuentes de error ya identificadas en la
Sección~\ref{sec:analisis_exp}.

Las barras de error $\pm\qty{0,88}{\milli\newton\per\meter}$ son
comparables a la dispersión observada entre puntos, lo que confirma que
la variabilidad experimental es consistente con las incertidumbres
instrumentales declaradas. La línea del valor bibliográfico
$\gamma_{\text{ref}} = \qty{72,8}{\milli\newton\per\meter}$ queda
notoriamente por encima de todos los puntos experimentales, evidenciando
una discrepancia sistemática que se discute en la
Sección~\ref{sec:conclusiones}.